\scanpage{250}
\mspage{196}

espérer s'en tirer avec \deleted{la} variante affaiblie de $W(4\mathrm b)$, obtenue en se bornant dans l'énoncé de $W(4\mathrm b)$ (p. 171) au cas où $X,Y$ sont dans $\underline{\mathrm{Perf}}_{W_{\unsure{\omega|w}}''}S$. Cette fois, ça suffit en effet à impliquer que le foncteur
\[
\lambda:(\underline{\mathrm{Fib}}_W S)(W_S^f)^{-1}\to(\underline{\mathrm{Fib}}_W S)(W_S)^{-1}
\]
est une équivalence de catégories, puisque on peut le ``calculer'' aussi bien en utilisant $\underline{\mathrm{Perf}}_{W_\omega''}S$ au lieu de $\underline{\mathrm{Fib}}_W$, or \deleted{\unsure{cette fois}} l'hyp. $W(4\ \mathrm b(W_\omega''))$ implique que sur $\underline{\mathrm{Perf}}_{W_\omega''}S$ les deux localiseurs \inserted{$W_S^f,W_S$} sont \uline{égaux}.

Il faut une conclusion pratique utilisable de toutes ces perplexités et observations. Essayons

\uline{Proposition} Supposons que l'on soit sous l'une des conditions suivantes

$1^\circ)$ L'axiome $W(4\mathrm b)$ est satisfait dans le cas de catégories $X,Y$ \deleted{\unsure{qui se}} qui sont dans $\underline{\mathrm{Fib\,sc}}_{\circ W}S$. (Cette condition \inserted{sur $W$} est en fait \unsure{autoduale|autodual.}.)
% slanted
\margin{mais c'est \uline{équivalent} à $W(4\mathrm b)$ ! cf $\overline{\underline{\mathrm{XVI}}}$ 4.4. loc. cit.}

$2^\circ)$ L'axiome $W(4\mathrm b)$ est satisfait dans le cas de catégories $X,Y$ qui sont dans $\underline{\mathrm{Perf}}_{W_\omega''}S$, où $W_\omega''$ est la plus petite ens. $\Sigma\subset\mathrm{Fl}(\mathrm{Cat})$ satisfaisant les conditions $W(1,2\ \mathrm{bis},3,6\ \mathrm{bis})$. (On peut aussi se contenter de l'a\deleted{\unsure{x}} condition, plus forte mais plus \unsure{compréhensible}, que $W(4\mathrm b)$ est OK si $X,Y$ sont dans $\underline{\mathrm{Perf}}_W S$ !) De plus, on suppose que $W$ satisfait $W(6\ \mathrm{bis})$.
% slanted
\margin{mais cela implique aussi $W(4\mathrm b)$ cf loc. cit.}

Sous ces conditions, le foncteur \inserted{\unsure{can.}}
\[
\mathrm{HOT}_W^{\unsure{\mathrm{lc}|\mathrm{le}}}(S)\to\mathrm{HOT}_W^{\unsure{\mathrm{lc}!|\mathrm{le}!}}(S)
\]
est une équiv. de cat., (cf diagramme p. \unsure{178|170}), donc $(\unsure{\rho|g},i);(\unsure{\sigma|\rho},j)$ \deleted{sont} \inserted{\unsure{peuvent} être regardés comme} des couples de foncteurs entre deux catégories (cf \unsure{(15)} p. 181). \uline{Ce sont des couples de foncteurs adjoints}, donc $i,j$ sont pl. fidèles, et on a les formules d'adjonction \unsure{(16)}, page 182.
% archive sheet number: 249
